\section{问题分析}

\subsection{几何与降维}

药材为圆柱，长 $L_0=0.25$ m、初始半径 $R_0=0.02$ m。
题目所有输出只到``距中心 2 cm''，即最大距离恰为初始半径，说明\textbf{径向}是主方向。
端面与侧面的面积比约为 $R_0/L_0=0.08$；端面方向与径向的扩散时间尺度比为
\begin{equation}
  \frac{t_z}{t_r}\sim\left(\frac{L_0/2}{R_0}\right)^{2}
  =\left(\frac{0.125}{0.02}\right)^{2}\approx39.06 .
\end{equation}
即端面方向的时间尺度比径向\textbf{大近 40 倍}，故本文采用
\textbf{一维轴对称径向模型}，并把端面效应作为独立对照定量验证（\S\ref{sec:endface}）。

\begin{figure}[H]
  \centering
  \includegraphics[width=0.98\textwidth]{G1_geometry.png}
  \caption{几何、边界与求解路线。(a) 横截面与侧面边界；(b) 侧视图，长径比 12.5；
           (c) 四问递进与 M0/M1/M2 分层}
  \label{fig:G1}
\end{figure}

\subsection{变量定义}

严格区分两个同为 kg/kg 但不能相减的量：
\begin{equation}
  C=\frac{m_w}{m_d}\quad(\text{药材干基含水率}),\qquad
  Y=\frac{m_v}{m_{da}}\quad(\text{空气含湿量}).
\end{equation}
本文另引入 $\Cinf$ 表示\textbf{已明确基准}的等效环境平衡含水率，
即附件1 的水分浓度列按题设等效约定读入的结果。

\subsection{技术路线}

技术路线如图 \ref{fig:G1b} 所示：从问题重述出发，先明确物理机制与变量定义，
再建立偏微分方程模型，采用有限体积法（FVM）与隐式时间推进求解，
依次回答四问，最后通过解析基准、收敛性、守恒、退化、灵敏度与稳健性六类检验验证模型。

\begin{figure}[H]
  \centering
  \includegraphics[width=0.86\textwidth]{G1b_flowchart.png}
  \caption{技术路线流程图}
  \label{fig:G1b}
\end{figure}

\subsection{四问递进}

\begin{itemize}[leftmargin=1.6em, itemsep=2pt]
  \item \textbf{问题 1}：预热平衡阶段，采用附录2 常物性，$D$ 不含 $T$；
  \item \textbf{问题 2}：整个烘干过程，采用附录3 变物性，$D(C,T)$ 双向耦合；
  \item \textbf{问题 3}：在问题2 模型上，确定全域最大含水率 $\Cmax$ 首次低于
        $0.15$ kg/kg 的时刻；
  \item \textbf{问题 4}：引入材料坐标与均匀径向仿射收缩，采用附录4 物性。
\end{itemize}

\subsection{模型分层}

\begin{itemize}[leftmargin=1.6em, itemsep=2pt]
  \item \textbf{M0}：题设基线，按附录与给定边界回答四问，\textbf{必做}，
        是四问的正式答案；
  \item \textbf{M1}：可核查扩展，潜热、辐射、端面、结壳、参数情景，
        \textbf{条件充分时做}，\textbf{不修订 M0 答案}；
  \item \textbf{M2}：工程推广，烘房系统、品质、结构，\textbf{只给路线与数据需求}。
\end{itemize}

% ==========================================================================
\subsection{本文的创新点}
\label{sec:innovation}

\begin{enumerate}[leftmargin=2.2em, itemsep=6pt]
  \item \textbf{收缩体的材料坐标处理：把移动边界化成固定域，且不引入伪对流。}

    常见处理有两种，都有代价：直接在物理坐标下重划网格，则每步都要重算网格与插值，
    离散引入的数值扩散会\textbf{伪装成物理扩散}；把坐标换成 $R_0$ 归一化却不改方程，
    则\textbf{悄悄丢掉}收缩带来的通量集中效应。

    本文改用材料坐标 $\xi=r/R(t)$，在均匀径向仿射收缩假设下（$v_r=\xi\dot R$）
    严格推得 $\dot R$ 对流项\textbf{恒为零}——这一步是\textbf{可证的，不是近似}；
    且扩散项 $\propto1/R^{2}$、表面项 $\propto1/R$，\textbf{二者不可合并}
    （很多文献把两者都写成 $1/R$ 或都写成 $1/R^{2}$，
    会在 $R$ 减半时产生百分之几十的系统误差）。
    因此问题4 与问题2/3 \textbf{共用同一套 FVM 内核}，
    $R(t)\equiv R_0$ 时逐点退化，这给了问题4 一个内建回归（已在 \S\ref{sec:degrade} 通过）。

  \item \textbf{``可忽略性必须绑定判据口径''——一个可复用的诊断范式。}

    端面效应在文献里常被引为``$3.5\%\sim6.0\%$，不可忽略''。
    本文建立二维轴对称有限圆柱模型做对照，发现：

    \begin{table}[H]
      \centering
      \small
      \begin{tabular}{@{}lccl@{}}
        \toprule
        用什么口径衡量端面效应 & 实测大小 & 结论 \\
        \midrule
        \textbf{全域最大含水率} $\Cmax(t)$（题面``各处''之意，本文判据） &
          $<10^{-5}$ & \textbf{可忽略} \\
        体积平均含水率 $\bar C$ & $-2.2\%\sim-6.2\%$ &
          与文献 $3.5\%\sim6.0\%$ \textbf{吻合} \\
        \bottomrule
      \end{tabular}
    \end{table}

    \textbf{同一个物理效应，两种口径差 3 个数量级。}
    本文的结论不是``端面效应可忽略''，而是：
    \textbf{任何``可忽略 / 不可忽略''的判断，都必须同时给出它绑定在哪个判据口径上；
    脱离口径讨论量级是无意义的。} 文献之间的分歧往往不是物理分歧，而是口径分歧。

  \item \textbf{问题1 的结构性单向耦合}：附录2 的 $D=7\times10^{-9}e^{-0.89/C}$
    \textbf{不含温度}，故问题1 可严格拆成``水分场独立求解 $\to$ 温度场代入已知
    $C(r,t)$ 求解''，不必双向迭代。本文据此把问题1 做成串行两步，
    并\textbf{用双向耦合求解器交叉验证}。

  \item \textbf{A/B/C 顺序差分}：把``物性公式''与``几何收缩''两个效应拆开。
    问题4 相对问题3 同时换了两样东西，若只报端点差 $C-A$，无法归因。
    本文按 A（附录3$+$固定）$\to$ B（附录4$+$固定）$\to$ C（附录4$+$收缩）
    分别报告 $\Delta_{\text{物性}}=B-A$ 与 $\Delta_{\text{几何}}=C-B$
    （见 \S\ref{sec:abc}）；并明确这\textbf{不是``唯一、独立的因果贡献分解''}。

  \item \textbf{三类误差分开报告，不把定位精度冒充整体精度}：
    $\tstar$ 的误差分解为事件定位 $0.5$ ms / 时间积分 $1.86$ s /
    \textbf{空间离散 $17.4$ s（主导）}，故 $\tstar$ 只报到 2 位小数，
    \textbf{不写 57.4889 h}——后者会让人误以为第 4 位有效。

  \item \textbf{几何—密度一致性诊断：条件性结论，不越界下判}。
    附录3 的 $\rho(C)=650+128C$ 与``干物质守恒''不兼容：
    $\rho_d=\rho(C)/(1+C)$ 在 $C=0$ 时为 $650$、在 $C=2.55$ 时为 $275$ kg/m$^3$。
    本文把它做成\textbf{条件性诊断}并明确写出：若 $\rho$ 是\textbf{热学等效参数}，
    则不应拿它做严格干质量检验；本文\textbf{不下``数据有误''的结论，
    也不为通过诊断而调整题设参数}。

  \item \textbf{显式精度口径声明：把``必须假定的东西''摊开}。
    附件1 只覆盖 0--14400 s，而问题2--4 要跑 2--3 天，
    恒温段边界\textbf{必须假定}。本文不隐藏这一点，而是给出 4 种同样合理的取法下
    $\tstar$ 的取值表（跨度 $0.30$ h $=1090$ s），
    并指出该跨度\textbf{比本文全部数值离散误差（约 19 s）大 57 倍}。

  \item \textbf{被省略因素的判据化取舍}：本文识别出 13 条文献中提及、
    但未进入四问答案的因素，给出\textbf{互斥且穷尽}的三类判据
    （量级 / 数据 / 重复计算），使每条因素都有明确归属，
    杜绝``文献有但我没做''式的无判据省略（见 \S\ref{sec:omitted}）。
\end{enumerate}
